\section{New workloads}
\label{sec:newworkloads}

\begin{question}
  \textbf{What do workloads beyond the campaign's two get from these nodes?} Kimi K3 inference first: what two nodes serve with vLLM, at what latency, how the cost grows with the context, and whether the answers stay accurate. Then Kimi K3 pretraining: whether one copy of the full model trains on the usable nodes, at what speed, and with how much memory to spare.
\end{question}

\paragraph{Kimi K3 inference, answer.} Two B200 nodes serve Kimi K3 with vLLM at up to \TriKimiPeakTok{} output tokens/s (\TriKimiPeakGpu{} per GPU) with \TriKimiPeakConc{} requests of 1k tokens in and 256 out in flight, \TriKimiSixtyFourTok{} with 64, and \TriKimiOneTok{} to a single request, one token every \TriKimiOneTpot\,ms. Throughput stops growing at \TriKimiPeakConc{} requests, the server's limit on sequences in flight, while its KV cache is about half full.

\paragraph{Setup.} Kimi K3 (\texttt{moonshotai/Kimi-K3}, about 2.8T parameters; a mixture of experts whose routed experts ship in MXFP4, the rest in BF16; text only here) does not fit on one node. It runs on \TriKimiNodes{} nodes, \TriKimiGpus{} GPUs: vLLM \TriKimiVllm{}, tensor parallelism over the 8 GPUs of each node and one pipeline stage per node, over Ray (\TriKimiLayout). vLLM runs the MXFP4 experts with activations quantized to MXFP8 (its \texttt{\TriKimiMoeBackend} backend), where the checkpoint expects BF16 activations, and FlashInfer's kernel autotuning is on. The server holds at most \TriKimiMaxSeqs{} sequences, up to \TriKimiMaxLen{} tokens each, and prefills in chunks of \TriKimiChunk{} tokens; prefix caching is off. Against it, the client runs the MI355X study's \TriKimiPasses{} passes unchanged, with random tokens and the end of sequence ignored: 1 to 1,024 concurrent requests of 1k tokens in and 256 out; inputs of 1k to 128k tokens, 128 out, at 1 and at 16 concurrent requests; and long decodes of 4k tokens. Then an accuracy pass: lm-eval through the server, temperature 0, up to 8,192 generated tokens. A throwaway warm-up pass runs first, so that kernels that compile on their first batches do so outside the measured passes. One run (\cref{tab:kimi_runs}), on \texttt{\TriKimiMainHost} and its neighbour; \texttt{\TriKimiMainHost} was held at 200\,W until November~6 (\cref{sec:machine}), and its GPUs averaged up to \TriKimiMainGpuW\,W each here, so that limit is gone. It is also one of the five nodes whose GPU boards run 1-node Qwen a little slower; any effect on this run is not measured.

\begin{table}[h]
  \caption{Kimi K3 inference runs in this report.}
  \label{tab:kimi_runs}
  \centering\small
  \IfFileExists{generated/tables/kimi_runs.tex}{\input{generated/tables/kimi_runs}}{\textit{No data yet.}}
\end{table}

\paragraph{What was measured.} All \TriKimiPasses{} passes ran: \TriKimiRequests{} requests, \TriKimiFailed{} failed (\cref{tab:kimi_passes,fig:kimi_inference}). \textbf{Concurrency.} Output grows from \TriKimiOneTok{} tokens/s for one request to \TriKimiSixtyFourTok{} at 64 and \TriKimiPeakTok{} at \TriKimiPeakConc, while the median time per output token grows from \TriKimiOneTpot{} to \TriKimiSixtyFourTpot{} and \TriKimiPeakTpot\,ms. At \TriKimiTopConc{} requests, output is \TriKimiTopPct\,\% lower than at \TriKimiPeakConc, and the median time to first token is \TriKimiTopTtft\,s: the server's log shows \TriKimiMaxSeqs{} requests running and the rest waiting. \textbf{Context.} For one request at a time, the median time to first token grows from \TriKimiCtxOneTtft\,ms at 1k tokens to \TriKimiCtxTtft\,ms at \TriKimiCtxTopK k, about in proportion to the prompt from 32k on, and the time per output token from \TriKimiCtxOneTpot{} to \TriKimiCtxTpot\,ms (+\TriKimiCtxTpotPct\,\%). With 16 at a time, each request waits for the others' prompts: \TriKimiCtxSixteenTtft\,s to the first token at \TriKimiCtxTopK k. From 128 concurrent requests up, and with 16 prompts of 8k tokens or more, the 99th percentile of the time between tokens rises to between 0.45 and 2\,s. \textbf{Accuracy} (\%, $\pm$ standard error): GSM8K \TriKimiGsm{} $\pm$ \TriKimiGsmErr, ARC-Challenge \TriKimiArc{} $\pm$ \TriKimiArcErr, GPQA diamond \TriKimiGpqa{} $\pm$ \TriKimiGpqaErr, and MMLU-Pro \TriKimiMmlu{} $\pm$ \TriKimiMmluErr{} on \TriKimiMmluN{} of its 12,032 questions. GSM8K and ARC-Challenge sit near their ceiling: they would catch a broken setup, not a small loss.

\begin{figure}[h]
  \centering\fig{kimi_inference}
  \caption{Kimi K3 inference on \TriKimiNodes{} B200 nodes (vLLM, \TriKimiLayout). Top: output tokens/s per GPU, and time per output token (median and 99th percentile), vs concurrent requests of 1k tokens in and 256 out; gray line: the server's limit of \TriKimiMaxSeqs{} sequences. Bottom: median time to first token and time per output token vs input length, 128 tokens out, at 1 and 16 concurrent requests. Log scales except top left.}
  \label{fig:kimi_inference}
\end{figure}

\begin{table}[h]
  \caption{Kimi K3 inference: the study's passes on \TriKimiNodes{} B200 nodes, random tokens, end of sequence ignored; the warm-up pass is left out. Per GPU: divided by \TriKimiGpus; per node: by \TriKimiNodes. Power is an estimate: the mean GPU power of the main node over each pass, times \TriKimiNodes, as the second node was not sampled; joules per output token follow from it.}
  \label{tab:kimi_passes}
  \centering\scriptsize
  \begin{adjustbox}{max width=\linewidth}
  \IfFileExists{generated/tables/kimi_passes.tex}{\input{generated/tables/kimi_passes}}{\textit{No data yet.}}
  \end{adjustbox}
\end{table}

\paragraph{Inferred.} The plateau from \TriKimiPeakConc{} requests on comes from the limit of \TriKimiMaxSeqs{} sequences, not from memory: vLLM sized the KV cache for \TriKimiKvTokens{} million tokens, its capacity, and with \TriKimiMaxSeqs{} sequences running it was at most \TriKimiKvUsePct\,\% full. If its use grows in proportion to the sequences, about \TriKimiKvFillSeqs{} would fill it; that was not run. The long tail of the time between tokens is likely decode steps waiting behind prefill chunks of \TriKimiChunk{} tokens. \textbf{Power and energy are estimates}, the main node's GPU power times \TriKimiNodes: about \TriKimiPowerOne\,W and \TriKimiJOne\,J per output token for one request, \TriKimiPowerSixtyFour\,W and \TriKimiJSixtyFour\,J at 64, and \TriKimiPowerPeak\,W and \TriKimiJPeak\,J at \TriKimiPeakConc{} (all est.). The two pipeline stages need not draw the same power, so the estimate can be off either way.

\paragraph{Missing.} The comparison with MI355X, pass for pass: its numbers are not in yet. The 4-node step, not run yet. The second node's power. A repeat: the study ran once.

\paragraph{Kimi K3 pretraining, answer.} Kimi K3 at full size, \TriKtLayers{} layers and \TriKtParams T parameters from random weights, trains as one copy on \TriKtNodes{} B200 nodes (\TriKtGpus{} GPUs, TorchTitan): \TriKtStep\,s per step of \TriKtTokens{} tokens (median of \TriKtMeasured{} measured steps), so \TriKtTok{} tokens/s per GPU and an MFU of \TriKtMfu\,\%, counting \TriKtFlops{} FLOPs per token as TorchTitan does. The fullest GPU reserved \TriKtMemWorst{} GiB of the \TriKtMemLimit{} GiB torch can use, \TriKtHeadroom{} GiB to spare. It fits with two deviations from the campaign's recipe: gradients are reduced in BF16, not FP32, and there are no FP32 master weights (parameters and optimizer state in BF16). In bring-up, FP32 master weights ran out of memory already on 1 and 2 nodes, and with FP32 reduction two layers per stage did not fit.

\paragraph{Setup.} Kimi K3's architecture (\texttt{moonshotai/Kimi-K3}, text only, no multi-token prediction): 896 routed experts, 16 of them and 2 shared ones active per token, linear attention (KDA) in 3 of 4 layers and MLA in the others, and attention residuals over blocks of 12 layers; \TriKtActive B parameters are active per token. TorchTitan's own Kimi K3 model, at the commit the campaign's Qwen pack pins. One copy spans the machine: a pipeline of \TriKtPp{} stages, one node each, two layers per stage except the last three (one each), with the node's 8 GPUs sharing the stage's experts (expert parallelism \TriKtEp) and its other weights (FSDP). Each of the 8 takes \TriKtMicro{} micro-batches of one \TriKtSeq-token FineWeb sequence per step, in a 1F1B schedule: \TriKtTokens{} tokens per step. AdamW, full activation recomputation, 20 steps of which the first 3 are warm-up. The \TriKtNodes{} nodes are all the usable ones (\cref{sec:machine}). \TriKtRuns{} runs at this scale (\cref{tab:kimi_train_runs}): the full-size one, one with \TriKtSmallLayers{} layers, and \TriKtFailed{} that stopped early.

\begin{table}[h]
  \caption{Kimi K3 pretraining runs on 48 nodes (PP48 $\times$ EP8, 384 GPUs). Step times on milabench's timer, steps 4-20; tokens/s per GPU and MFU at the median step, MFU with TorchTitan's FLOPs per token. Step 1: minutes from the job's start, from the log of node 0 (``--'' where the run's data file holds none). Why runs stopped: the experiment coordinator's and the implementers' diagnoses (\texttt{RUNS.md}), not measured here.}
  \label{tab:kimi_train_runs}
  \centering\scriptsize
  \begin{adjustbox}{max width=\linewidth}
  \IfFileExists{generated/tables/kimi_train_runs.tex}{\input{generated/tables/kimi_train_runs}}{\textit{No data yet.}}
  \end{adjustbox}
\end{table}

\paragraph{What was measured} (full size, \cref{fig:kimi_training}). \textbf{The first two steps.} Step 1 ended \TriKtStepOne{} min after the job started: before it, the pipeline infers each stage's shapes, one stage after another. Step 2 took \TriKtStepTwo\,s: its first backward pass reaches the stages one after another, and each stage pauses once as it does (GPU 0's utilization drops for some seconds), about \TriKtSweepHi\,s per stage from stage \TriKtSweepFirst{} down to stage \TriKtSweepSplit{} and \TriKtSweepLo\,s below it, \TriKtSweepTotal\,s in all. What each stage does in that pause is not identified. \textbf{Steps \TriKtMeasuredFrom-20.} Median \TriKtStep\,s, 90th percentile \TriKtStepPninety, slowest \TriKtStepMax; \TriKtTokLo{} to \TriKtTokHi{} tokens/s per GPU step by step. The step time drifted, \TriKtStepFirst\,s at step \TriKtMeasuredFrom, \TriKtStepMax{} at its slowest and \TriKtStepLast{} at step 20. TorchTitan's own step intervals agree with milabench's within a few milliseconds. With the FLOPs of 6 $\times$ the active parameters alone, without attention's, MFU would be \TriKtMfuSixN\,\%. \textbf{Memory} (peak per GPU, the highest of a stage's 8): stage \TriKtMemWorstStage{} (\texttt{\TriKtMemWorstHost}) reserved the most, \TriKtMemWorst{} GiB (\TriKtMemWorstAlloc{} allocated). Stages 6, 12, \dots, 42 reserved \TriKtMemEdgeLo-\TriKtMemEdgeHi{} GiB, the other two-layer stages \TriKtMemMidLo-\TriKtMemMidHi; stage 0, with the embedding, \TriKtMemZero{} (\TriKtMemZeroAlloc{} allocated); the last stage, one layer and the output, \TriKtMemLast{} (\TriKtMemLastAlloc). No allocator failure, no error. \textbf{Loss} from \TriKtLossFirst{} at step 1 to \TriKtLossLast{} at step 20, with the same gradient norm logged by all 48 stages at every step.

\begin{figure}[h]
  \centering\fig{kimi_training}
  \caption{Kimi K3 pretraining on 48 B200 nodes. Left: step time of the measured steps (milabench's timer). Middle: peak memory per GPU by pipeline stage, the highest of each stage's 8 GPUs; solid: reserved, dashed: allocated (full size only); red: what torch can use. Right: training loss. The two runs differ in layers, gradient reduction and the last stage's loss handling (text), so they do not compare one to one.}
  \label{fig:kimi_training}
\end{figure}

\paragraph{Inferred.} Stages 6, 12, \dots, 42 hold the first layer of each 12-layer block of attention residuals (from TorchTitan's model code, not measured), which is why they reserve more than their neighbours. The step is longer than estimated before the run (65-80\,s per step and 8.0-9.9\,\% MFU, est.). Two parts of the step are measured: fetching a step's data takes \TriKtFetch\,s on each rank, inside the step, and GPU 0 of every stage draws under 350\,W for \TriKtLowLo-\TriKtLowHi\,\% of the measured steps, about the share 1F1B leaves idle with 48 stages and 128 micro-batches, \TriKtBubblePct\,\% (computed, $(p-1)/(m+p-1)$). The rest of the gap to the estimate, and the drift, are not determined.

\paragraph{What it took at 48 stages.} \TriKtFailed{} earlier runs at this scale stopped early, each on one way a deep pipeline breaks (\cref{tab:kimi_train_runs}). \textbf{Timeouts.} The first step infers the stages' shapes one stage at a time, and step 2's first backward pass crosses the stages one at a time too, so each takes many minutes at 48 stages. The pipeline's own communication groups kept a 10-minute timeout that the init timeout did not reach, and timed out in the first step; and after step 1 TorchTitan lowers every group to its train timeout, 3 minutes by the pack's default, which step 2 outlasted. Both timeouts were raised above the warm-up (60 and 30 minutes). \textbf{Memory.} The last stage holds one layer and the output, but torch's 1F1B schedule kept every micro-batch's loss until the step ended, and with it the micro-batch's graph and the inputs the stage received; with 128 micro-batches it ran out of memory (found on a 4-node stand-in of that stage; the 48-node run recorded no breakdown). Detaching each loss once its backward has run keeps it at \TriKtMemLast{} GiB. \textbf{Monitoring.} One run was cancelled after step 1 under a rule, no new step line for 5 minutes, that could not work with a log line every 10 steps; it was probably still training.

\paragraph{Fewer layers.} The \TriKtSmallLayers-layer run, one layer per stage, differs from the full-size one in three ways: half the layers, FP32 gradient reduction, and no loss detach. So the two are not a clean comparison. Without the detach, its last stage reserved \TriKtSmallMemLast{} GiB, close to the limit, against \TriKtSmallMemZero{} on stage 0; its step 2 took \TriKtSmallStepTwo\,s, about \TriKtSmallSweep\,s per stage. Its median step, \TriKtSmallStep\,s (\TriKtSmallStepLo-\TriKtSmallStepHi), is longer than the full-size run's, with MFU \TriKtSmallMfu\,\%; its last stage logged 20 allocator warnings (``memory mapping failed''), which may explain the slow steps (inferred, not measured).

\paragraph{Missing.} Twenty steps only: nothing on convergence, or on stability over a long run. The cause of the drift, and of the gap to the estimate. A repeat: one full-size run. 52 nodes, with the four nodes still out. FP32 master weights and FP32 gradient reduction, which did not fit in bring-up. The share of the step spent in the experts' all-to-all.

\paragraph{Still to come.} A block for GLM-5.3 RL.
